\documentclass[10pt,a4paper]{amsart}
\usepackage[margin=19mm]{geometry}
\usepackage{amsmath,amssymb,mathtools}
\usepackage[scr=rsfs,cal=euler]{mathalfa}
\usepackage{xfrac}
\usepackage[unicode,hidelinks,bookmarks=false]{hyperref}
\newcommand{\ie}{i.e.\ }
\newcommand{\cf}{cf.\ }
\newcommand{\resp}{respectively\ }
\newcommand{\Subsection}{\subsection}
\providecommand{\nobreakdash}{\nobreak\hbox{-}}
\setlength{\emergencystretch}{2em}
\multlinegap0pt
\usepackage{bm}
\usepackage{bbm}
\usepackage{dsfont}
\usepackage{xspace}
\usepackage{cancel}
\usepackage{mathtools}
\usepackage{amsthm}
\makeatletter
\@ifundefined{thm}{\newtheorem{thm}{Theorem}}{}
\@ifundefined{cor}{\newtheorem{cor}[thm]{Corollary}}{}
\@ifundefined{lem}{\newtheorem{lem}[thm]{Lemma}}{}
\@ifundefined{rem}{\newtheorem{rem}[thm]{Remark}}{}
\makeatother
\newcommand{\R}{\mathbb{R}}
\title[Critical quasi-linear Schr\"{o}dinger system with $p$-Laplac]{Critical quasi-linear Schr\"{o}dinger system with $p$-Laplacian}
\author{Neng cheng, Wei Dai, Zhao Liu}
\date{Source: arXiv:2606.10630v1 (CC BY 4.0)}
\begin{document}
\maketitle

\noindent\emph{Source and licence.} Critical quasi-linear Schr\"{o}dinger system with $p$-Laplacian. Version arXiv:2606.10630v1 (CC BY 4.0), distributed under the Creative Commons Attribution 4.0 International licence, as recorded in the arXiv licence metadata for this exact version and shown on the abstract page https://arxiv.org/abs/2606.10630v1. Full rights notice: licenses/arxiv-2606.10630-CC-BY-4.0.txt.\par\smallskip

\noindent\textit{Editorial scope.} Counted unit: the Lem whose source label is \texttt{lm.4.2} in the article (source0.tex lines 1379--1394, source label \texttt{lm.4.2}), delivered together with the article's own complete original proof of it (source0.tex lines 1395--1427, one \texttt{proof} environment of 33 lines). No mathematics is written, repaired or re-derived: statement and proof are the article's own text line for line. Required context retained uncounted: eq1.1 (lines 126--133); th2.1- (lines 206--215); re2333 (lines 338--344); eq10.25.2 (lines 340--342); re2334 (lines 352--354); lm.4.1 (lines 1181--1194); eq1204 (lines 1324--1332); re.4.1 (lines 1372--1376). Every definition, display and label of those lines is retained unchanged. Omitted from this package and not redistributed: 1--125, 134--205, 216--337, 343--351, 355--1180, 1195--1323, 1333--1371, 1377--1378, 1428--1767. Nothing inside a retained excerpt is omitted or altered: every retained body is delivered exactly as the article's lines are written, so no omission receipt of any kind accompanies this package. Citations: the retained excerpts carry no citation command at all, so no bibliography entry of the article is transcribed and no reference is redistributed. Reference adaptation: none was needed and none was made; every reference inside the delivered unit resolves inside it, so no unresolved reference or citation is produced. Printed slips kept literal: no printed word, symbol, hypothesis or number of the counted statement or of its proof was altered. Layout and class: the article's own class is \texttt{amsart} with packages including graphicx, amsmath, amsfonts, amssymb, setspace, datetime, hyperref, geometry; the wrapper is the lane's pinned \texttt{amsart} editorial wrapper at 10pt on A4, which restates the theorem-like environments the retained text needs and adds the article's own macro definitions that the retained text needs (\texttt{\textbackslash R}), with extra wrapper packages (bm, bbm, dsfont, xspace, cancel, mathtools). The delivered numbering is the unit's own. Assets: no retained span contains a figure environment, a graphics-inclusion command, a PDF-page inclusion, a table or a TikZ picture; no image file is copied into this package, the package declares no asset dependency and no image file is redistributed with it. Licence: version arXiv:2606.10630v1 of this article is distributed under the Creative Commons Attribution 4.0 International licence, as recorded in the arXiv licence metadata for this exact version and shown on the abstract page https://arxiv.org/abs/2606.10630v1. A full rights notice is delivered at licenses/arxiv-2606.10630-CC-BY-4.0.txt. Characters: every retained excerpt is pure ASCII, so the delivered engine, XeLaTeX with its default Latin Modern text font, reports no missing character.
\begin{align}\label{eq1.1}
\left\{ \begin{array}{ll} \displaystyle
-\Delta_p u = u^{\alpha}v^{\beta}  \, \ \ \ \ \ \mbox{in}\,\ \ \R^N, \\ \\
-\Delta_p v = u^{\beta}v^{\alpha}  \,\  \ \ \ \ \mbox{in}\,\ \ \R^N, \\ \\
u, v \in D^{1,p}(\R^N),\,\,\,u, v>0 \quad \mbox{in}\,\, \ \ \R^N,
\end{array}
\right.\hspace{1cm}
\end{align}

\begin{thm}
[Regularity and sharp asymptotic estimates]\label{th2.1-}
Assume $1<p<N$, and let $(u,v)$ be a pair of positive $D^{1,p}(\R^{N})$-weak solutions to the $D^{1,p}(\mathbb{R}^{N})$-critical quasi-linear Schr\"{o}dinger system \eqref{eq1.1}. Then $u, v \in C_{loc}^{1,\theta} (\R^N) \cap L^\infty(\R^N)$. Moreover, there exist $c_0, C_0,c'_0, C'_0, R_0 > 0$ such that
\begin{equation}\label{eq0806}
  \frac{c_0}{1+|x|^\frac{N-p}{p-1}} \leq u(x) , v(x) \leq \frac{C_0}{1+|x|^\frac{N-p}{p-1}} \,\,\,\,\,\,\,\,\mbox{in}\,\,\ \R^N,
\end{equation}
\begin{align}\label{eq0806+}
   \frac{c'_0}{|x|^\frac{N-1}{p-1}} \leq |\nabla u(x)|, |\nabla v(x)|\leq \frac{C'_0}{|x|^\frac{N-1}{p-1}} \,\,\,\,\,\,\,\,\mbox{in}\,\,\ \R^N\setminus B_{R_0}(0).
\end{align}
\end{thm}

\begin{cor}\label{re2333}
Let $1<p<N$ and $u$ be a weak positive solution of \eqref{eq1.1}. Setting $\rho_u:=|\nabla u|^{p-2}, \rho_v:=|\nabla v|^{p-2}$, for any $\Omega \subset\subset\R^N$, we have
\begin{align}\label{eq10.25.2}
\int_{\Omega} \frac{1}{\rho_u^t} \frac{1}{|x-y|^\gamma}\mathrm{d}x < C \quad and \quad \int_{\Omega} \frac{1}{\rho_v^t} \frac{1}{|x-y|^\gamma}\mathrm{d}x < C
\end{align}
uniformly for any $y\in\Omega$, where $\max\{p-2,0\}\leq (p-2)t<p-1$ and $\gamma=0$ if $N=2$, $\gamma<N-2$ if $N \geq 3$.
\end{cor}

\begin{rem}\label{re2334}
Let $(u,v)$ be a pair of positive weak solutions of \eqref{eq1.1} and let $Z_u = \{x \in \R^N|\ |\nabla u(x)| = 0 \}$. Clearly, it follows from Theorem \ref{th2.1-} that the critical set $Z_u\subset B_{R_0}$ and $Z_u$ is a closed set. Moreover, \eqref{eq10.25.2} in Corollary \ref{re2333} implies that the Lebesgue measure of $Z_{u}$ is zero, i.e., $|Z_u|=0$. In the same argument, we have $|Z_v|=0$.
\end{rem}

\begin{lem}\label{lm.4.1}
Assume $N>2$, $1<p<2$ and $(u, v)$ is a pair of positive weak solutions of \eqref{eq1.1}. Let $R_0$ be the same as in Theorem \ref{th2.1-}.

{\bf $(i)$} There exists $R_1>R_0$, $R_1$ depending on $N,p,R_0$, such that if
$$\Sigma_\lambda^\nu \cap R_\lambda^\nu(B_{R_1}) = \emptyset,$$
then $u^{\nu}_{\lambda} \leq u$ and $v^{\nu}_{\lambda} \leq v$ in $\Sigma_\lambda^\nu$;


{\bf $(ii)$} Let $\Sigma_\lambda^\nu \cap R_\lambda^\nu(B_{R_1}) \neq \emptyset$ and let $R_2=R_2(R_1,\lambda,\nu)>2R_{1}$ be such that $\Sigma_\lambda^\nu \cap R_\lambda^\nu(B_{R_1}) \subset B(P_\lambda^\nu,R_2)$, where $P_\lambda^\nu$ is the projection of the origin on $T_\lambda^\nu$.
Then there exists $R>R_2$, depending on $R_2$, such that, for any $\overline{R}\geq R$, there exist positive numbers $\delta$ and $M$, depending $\overline{R},\lambda,p,N$ and the $L^\infty$--norms of $u,v,|\nabla u|,|\nabla v|$, with the property that whenever
$$ (supp(u^{\nu}_{\lambda}-u)^+\cup supp(v^{\nu}_{\lambda}-v)^+)\cap\Sigma^{\nu}_{\lambda}\cap B(P_\lambda^\nu,\overline{R})=A \cup B,$$
with $|A\cap B|=0$, $M_A<M$ and $|A|+|B|<\delta$, we have
$$ u^{\nu}_{\lambda} \leq u\,\,\,\,\mbox{and}\,\,\,\, v^{\nu}_{\lambda} \leq v \,\,\quad \mbox{in}\,\,\Sigma_\lambda^\nu.$$
\end{lem}

\begin{align}\label{eq1204}
&\quad\int_{\Sigma^{\nu}_\lambda} (|\nabla u|+|\nabla u^{\nu}_{\lambda}|)^{p-2} |\nabla(u^{\nu}_{\lambda}-u)^+|^2 \mathrm{d}x \\
&\leq \frac{C_0(2\overline{C}_1+\overline{C}_2)}{2C'(2C'_0)^{p-2}} \frac{1}{R^{\beta_1}}\left(\frac{2}{N+s-2}\right)^2 \int_{\Sigma_\lambda^\nu}(|\nabla u|+|\nabla u^{\nu}_{\lambda}|)^{p-2} |\nabla (u^{\nu}_{\lambda}-u)^+|^2 \mathrm{d}x \nonumber\\
& \quad +\frac{ C_0\overline{C}_2}{2C'(2C'_0)^{p-2}}\frac{1}{R^{\beta_1}}\left(\frac{2}{N+s-2}\right)^2\int_{\Sigma_\lambda^\nu} (|\nabla v|+|\nabla v^{\nu}_{\lambda}|)^{p-2} |\nabla(v^{\nu}_{\lambda}-v )^+|^2 \mathrm{d}x  \nonumber\\
&\quad +\frac{\overline{C}_1\|v^\nu_\lambda\|^\beta_\infty }{C'}C |A_\lambda^\nu|^\frac{1}{N} \left(  |A_\lambda^\nu|^\frac{1}{N} M_A^{2-p}+|B|^\frac{1}{N} M_K^{2-p}\right)  \nonumber\\
&\quad \quad \times \int_{\Sigma_\lambda} (|\nabla u|+|\nabla u^{\nu}_{\lambda}|)^{p-2} |\nabla(u^{\nu}_{\lambda}-u )^+|^2 \mathrm{d}x \nonumber\\
 &\quad +\frac{\overline{C}_2\|u^\nu_\lambda\|^\alpha_\infty }{2C'}C |A_\lambda^\nu|^\frac{1}{N} \left(|A_\lambda^\nu|^\frac{1}{N} M_A^{2-p}+ |B|^\frac{1}{N} M_K^{2-p} \right)  \nonumber\\
&\quad \quad \times \left(\int_{\Sigma_\lambda} (|\nabla v|+|\nabla v^{\nu}_{\lambda}|)^{p-2} |\nabla(v^{\nu}_{\lambda}-v )^+|^2 +(|\nabla u|+|\nabla u^{\nu}_{\lambda}|)^{p-2} |\nabla(u^{\nu}_{\lambda}-u )^+|^2\right)\mathrm{d}x. \nonumber
\end{align}

\begin{rem}\label{re.4.1}
Observe that, for given $\lambda_0$ and $\nu_0$, $\overline{R}$ can be chosen so large such that, for all $\lambda$ in a small neighbourhood of $\lambda_0$, say, $[ \lambda_0-\theta_0 ,\lambda_0 + \theta_0 ]$ with $\theta_{0}>0$ small, and for all $\nu$ in $\overline{I_{\theta_0}(\nu_0)}:=\{ \nu \mid |\nu|=1,|\nu-\nu_0|\leq \theta_0 \}$, we have
$$ \overline{R}>R=R(\lambda,\nu)\,\quad \mbox{and}\,\quad B(P_\lambda^\nu,R)\subset B(P_{\lambda_0}^{\nu_0},\overline{R}), $$
where $R=R(\lambda,\nu)$ is given by Lemma \ref{lm.4.1} $(ii)$. Furthermore, $\delta$ and $M$ can also be chosen uniformly for all these $\lambda$ and $\nu$ in a $\theta_0$--neighbourhood of $(\lambda_0,\nu_0)$, by replacing $|A_\lambda^\nu|$ with $|B(P_{\lambda_0}^{\nu_0},\overline{R})|$ on the R.H.S of \eqref{eq1204}.
\end{rem}

\begin{lem}\label{lm.4.2}
Assume $(u,v)$ is a pair of positive weak solutions of \eqref{eq1.1} for $1<p<2$. Let $\nu_0$ be any unit vector in $\R^N$ and $\lambda_0\in\overline{\Lambda}(\nu_0)$. Let $\overline{R}$, $\delta$ and $M$ be given by Remark \ref{re.4.1} corresponding to $\theta_0-$neighbourhood of $(\lambda_0,\nu_0)$.

If we have
\begin{equation}\label{assump}
  u>u_{\lambda_0}^{\nu_0},\quad\,\,\mbox{in} \,\,\,(\Sigma_{\lambda_0}^{\nu_0}\setminus Z_{\lambda_0}^{\nu_0}(u))  \cap  B(P_{\lambda_0}^{\nu_0},\overline{R}),
\end{equation}
\begin{equation}\label{assump-}
  v>v_{\lambda_0}^{\nu_0},\quad\,\,\mbox{in} \,\,\,(\Sigma_{\lambda_0}^{\nu_0}\setminus Z_{\lambda_0}^{\nu_0}(v))  \cap  B(P_{\lambda_0}^{\nu_0},\overline{R}),
\end{equation}
then there exists $0<\theta_1 < \theta_0$ such that
%{\bf $(i)$} $u\geq u_\lambda^{\nu_0}$ in $\Sigma_\lambda^{\nu_0} \,\,\forall\,\lambda\in(\lambda_0-\theta_1,\lambda_0+\theta_1)$,\\
%{\bf $(ii)$}
$$ u\geq u^{\nu}_{\lambda},\,\,v\geq v^{\nu}_{\lambda} \quad \mbox{in}\,\, \Sigma_\lambda^\nu,\,\,\,\,\quad \forall\,\,(\lambda,\nu)\in (\lambda_0-\theta_1,\lambda_0+\theta_1)\times I_{\theta_1}(\nu_0),$$
where $I_{\theta_1}(\nu_0):=\{ \nu \,|\, |\nu|=1, |\nu-\nu_0|<\theta_1 \}$.
\end{lem}

\begin{proof}
Let $K:=B(P_{\lambda_0}^{\nu_0},\overline{R})$, where $P_{\lambda_0}^{\nu_0}$ is the projection of the origin on $T_{\lambda_0}^{\nu_0}$, and let $\delta$ and $M$ be chosen as in Remark \ref{re.4.1} uniformly for all $(\lambda,\nu)$ in a $\theta_0$-neighbourhood of $(\lambda_0,\nu_0)$.

Since $\overline{R}\geq R>R_{2}>2R_{1}>2R_{0}$, $\Sigma_\lambda^\nu \cap R_\lambda^\nu(B_{R_1}) \subset B(P_\lambda^\nu,R_2) \subset B(P_\lambda^\nu,R)\subset B(P_{\lambda_0}^{\nu_0},\overline{R})$ and $\Sigma_\lambda^\nu \cap R_\lambda^\nu(B_{R_1}) \neq \emptyset$ imply $B_{R_0}\subset B(P_{\lambda_0}^{\nu_{0}},\overline{R})$, by the sharp asymptotic estimates in Theorem \ref{th2.1-}, we get
$$Z_u\cup Z_v\subset K.$$
On the other hand, by Remark \ref{re2334}, the Lebesgue measures of $Z_{u}$ and $Z_{v}$ are zero, i.e., $|Z_u|=|Z_v|=0$.  Thus we can choose a small neighborhood $B$ of $\partial\left(K\cap \Sigma_{\lambda_0}^{\nu_0}\right)$ with $|B|<\frac{\delta}{4}$ and another neighbourhood $O$ of $(Z_{\lambda_0}^{\nu_0}(u)\cup Z_{\lambda_0}^{\nu_0}(v)) \cap K$ with $|O|<\frac{\delta}{4}$ and $M_O^{\lambda_0,\nu_0}<\frac{M}{2}$, where
$$ M_O^{\lambda,\nu}=\max\{\sup_O (|\nabla u|+|\nabla u^{\nu}_{\lambda}|), \sup_O (|\nabla v|+|\nabla v^{\nu}_{\lambda}|)\} .$$

Note that $K_0=(\overline{K\cap\Sigma_{\lambda_0}^{\nu_0}}) \setminus (B\cup O)$ is a compact set, so by our assumption \eqref{assump}, one has, there exists some positive constant $m>0$ such that
$$ u-u_{\lambda_0}^{\nu_0}>m>0,\  v-v_{\lambda_0}^{\nu_0}>m>0 \quad\quad \,\mbox{in}\,\,\, K_0. $$
Then the continuity with respect to $\lambda$ and $\nu$ implies that, there exists $\theta_1=\theta_1(\lambda_0,\nu_0,K)<\theta_0$, such that
\begin{align*}
& u-u^{\nu}_{\lambda}> \frac{m}{2}>0,\ v-v^{\nu}_{\lambda}> \frac{m}{2}>0 \qquad \,\mbox{in}\,K_0, \\
& M_{O\cap \Sigma_\lambda^\nu}^{\lambda,\nu}< M, \\
& |B\cap B(P_\lambda^\nu,R) \cap \Sigma_\lambda^\nu |< \frac{\delta}{2},\\
& |O\cap B(P_\lambda^\nu,R) \cap \Sigma_\lambda^\nu |< \frac{\delta}{2},\\
& (B(P_\lambda^\nu,R)\cap \Sigma_\lambda^\nu)\subset (B(P_{\lambda_0}^{\nu_0},\overline{R})\cap \Sigma_{\lambda_0}^{\nu_0}) \cup B
\end{align*}
for all $(\lambda,\nu)\in (\lambda_0-\theta_1, \lambda_0+\theta_1)\times I_{\theta_1}(\nu_0)$, where $R=R(\lambda_0,\nu_0)$ is given by Lemma \ref{lm.4.1} $(ii)$. Then we have
\begin{align*}
     & \quad(supp(u^{\nu}_{\lambda}-u)^+\cup supp(v^{\nu}_{\lambda}-v)^+) \cap B(P_\lambda^\nu,R)\\
     &\subset (B\cup O) \cap B(P_\lambda^\nu,R) \\
     &= (B\cap B(P_\lambda^\nu,R)) \cup (O\cap B(P_\lambda^\nu,R))
\end{align*}
for all $(\lambda,\nu)\in (\lambda_0-\theta_1, \lambda_0+\theta_1)\times I_{\theta_1}(\nu_0)$, which implies that
\begin{align*}
    & \quad(supp(u^{\nu}_{\lambda}-u)^+\cup supp(v^{\nu}_{\lambda}-v)^+) \cap B(P_\lambda^\nu,R)\cap \Sigma_\lambda^\nu\\
    &\subset(B\cap B(P_\lambda^\nu,R)\cap \Sigma_\lambda^\nu) \cup (O\cap B(P_\lambda^\nu,R)\cap \Sigma_\lambda^\nu)
\end{align*}
satisfies the conditions of Lemma \ref{lm.4.1} $(ii)$. Therefore, we derive
$$ u\geq u^{\nu}_{\lambda},\ v\geq v^{\nu}_{\lambda} \quad \,\,\,\,\,\mbox{in}\,\, \Sigma_\lambda^\nu$$
for all $(\lambda,\nu)\in (\lambda_0-\theta_1,\lambda_0+\theta_1)\times I_{\theta_1}(\nu_0)$. This completes our proof of Lemma \ref{lm.4.2}.
\end{proof}

\end{document}
